\label{section:vector-register-model}
The scalable-vector extension provides 32 (:term:VECTOR.terms.vector_register|plural:), \texttt{V0}
through \texttt{V31}, and 16 (:term:VECTOR.terms.predicate_register|plural:), \texttt{P0} through
\texttt{P15}.  A (:term:VECTOR.terms.vector_register:) contains VLEN bits.  VLEN is an
implementation-selected power of two from 128 through 2048 bits and remains
stable until reset.  A (:term:VECTOR.terms.predicate_register:) contains one bit for each byte of a
(:term:VECTOR.terms.vector_register:); its packed image therefore occupies VLEN/64 bytes.  Reset
clears every vector and predicate bit.

The \texttt{VECTOR} discovery predicate and the \texttt{VECTOR\_PARAMETERS}
CPUID leaf report the extension and VLEN. Integer and raw-bit vector forms
require \texttt{VECTOR}. Floating-point vector operations are provided by the
separate \texttt{VECTORFP} extension.

\subsection{Element Types and Assembly Spelling}

A size-bearing VECTOR instruction has one suffix. The integer element
suffixes \texttt{B}, \texttt{W}, \texttt{L}, and \texttt{Q} select 1-, 2-,
4-, and 8-byte elements.

The three-bit \texttt{tzz} selector values zero through three select
\texttt{B/W/L/Q}; value four is invalid. The two-bit \texttt{zz} selector in
integer families assigns all four values to \texttt{B/W/L/Q}. A
(:term:base.terms.reserved:) selector does not name another data format.

The suffix is the source width for a width-changing instruction; the
destination width and signedness are part of the mnemonic. For example,
\texttt{VEXTZQ.B} zero-extends a byte into a quadword. Predicate
operations, (:ref:VECTOR.instructions.VLCNT:), (:ref:VECTOR.instructions.VLCADD:), (:ref:VECTOR.instructions.VGATHER:),
(:ref:VECTOR.instructions.VSCATTER:), (:ref:VECTOR.instructions.PLOOP:), (:ref:VECTOR.instructions.VMOV:), (:ref:VECTOR.instructions.VMOVZ:),
and the permute, slide, slice, and zip families use \texttt{B/W/L/Q};
\texttt{H/S/D} width-only aliases are not defined.
